\section{A lattice-theoretic lead and a sharp small-defect boundary}
\label{sec:completion-lead}
The classical completion theorem suggests studying the orthogonal
complement of a partial Hadamard matrix through arithmetic quadratic
forms. It is not an automatic solution: rational orthogonality,
integrality and sign entries are different requirements.

Let $X$ have $n-r$ orthogonal sign rows of length $n$, so
$XX^T=nI$. A real orthogonal completion $Y$ satisfies
\[
 Y^TY=nI-X^TX.
\]
The right side has rank $r$, diagonal entries $r$, integer entries,
and rigid congruences inherited from signs. Completing $X$ to a
Hadamard matrix asks for a factorization of this specific Gram matrix
by an $r\times n$ \emph{sign} matrix. A general real, rational, or
integer Gram factorization is a weaker problem.

\begin{proposition}[Eight missing rows can already be fatal]
\label{prop:eight-defect}
Let $K$ be a Hadamard matrix of order $m$ and let $s>1$ be odd.
Repeating every column $s$ times gives an $m\times ms$ partial
Hadamard matrix with no additional orthogonal sign row.
In particular there is an unextendible $4\times12$ example, with
eight missing rows, although Hadamard matrices of order $12$ exist.
\end{proposition}
\begin{proof}
The repeated matrix is $X=K\otimes(1,\ldots,1)$, and
$XX^T=msI_m$. For a putative new sign row, let $b_j$ be its sum
on the $j$th length-$s$ block. Orthogonality says $Kb=0$, so $b=0$
by invertibility of $K$. Each $b_j$ is a sum of an odd number of
signs and cannot be zero. Take $m=4,s=3$ for the example.
\end{proof}
This example is classical: the $4\times12$ pattern is displayed
in Verheiden's Section 4~\cite{Verheiden}. Our proof makes its block
obstruction explicit. Combined with his seven-row theorem, it shows
that universal completion based only on the number of missing rows
cannot be pushed from seven to eight.

The thesis also uses arithmetic results of Kneser and others to
restrict denominators of rational Gram completions. This is a
concrete route for further research: classify the residual Gram
forms and identify additional structure permitting sign factorizations
at rank eight and beyond. An $E_8$ explanation or an all-rank
completion theorem is not established here. In particular, the
unextendible example shows that any useful theorem must select or
modify the partial rows, rather than complete every possible choice.
